% \begin{table}[H]
% \centering
% \renewcommand{\arraystretch}{1.2}
% \begin{tabular}{@{}p{4cm} p{11cm}@{}}
% \toprule
% \textbf{Use case ID}       & UC-M8-01 \\ \midrule
% \textbf{Tên Use case}      & Đăng ký tài khoản \\ \midrule
% \textbf{Các tác nhân}      & Người dùng chung \\ \midrule
% \textbf{Mô tả}             & Tạo tài khoản mới cho người dùng\\ \midrule
% \textbf{Trigger}    & Người dùng bấm “Đăng ký”.\\ \midrule
% \textbf{Tiền điều kiện}    & Người dùng chưa đăng nhập và email/sđt dùng để đăng ký chưa tồn tại.
% \\ \midrule
% \textbf{Hậu điều kiện}     & Tài khoản mới được thêm vào hệ thống. \\ \midrule
% \textbf{Luồng sự kiện chính} & 
% \begin{tabular}[t]{@{}p{10.5cm}@{}}
% 1. Người dùng nhấn nút “Đăng ký”.\\
% 2. Hệ thống hiển thị trang có form đăng ký tài khoản.\\
% 3. Người dùng nhập thông tin để đăng ký tài khoản.\\
% 4. Hệ thống kiểm tra định dạng \& trùng lặp. \\
% 5. Hệ thống tạo tài khoản. 
% \end{tabular} \\ \midrule
% \textbf{Luồng ngoại lệ} &
% \begin{tabular}[t]{@{}p{10.5cm}@{}}
% 4a. Email/SĐT đã tồn tại → báo lỗi “Tài khoản đã tồn tại”. 
% 5a. Hệ thống gặp sự cố khi tạo tài khoản → báo lỗi “Lỗi kết nối”. 
% \end{tabular} \\
% \bottomrule
% \end{tabular}
% \caption{Mô tả Use Case: Đăng ký tài khoản}
% \label{usecase-8-01}
% \end{table}

% \begin{table}[H]
% \centering
% \renewcommand{\arraystretch}{1.2}
% \begin{tabular}{@{}p{4cm} p{11cm}@{}}
% \toprule
% \textbf{Use case ID}       & UC-M8-04 \\ \midrule
% \textbf{Tên Use case}      & Đăng nhập bằng tài khoản Google \\ \midrule
% \textbf{Các tác nhân}      & Người dùng chung \\ \midrule
% \textbf{Mô tả}             & Cho phép người dùng đăng nhập vào hệ thống qua Google.\\ \midrule
% \textbf{Trigger}    & Người dùng bấm “Đăng nhập với Google".\\ \midrule
% \textbf{Tiền điều kiện}    &  Người dùng đã có tài khoản trong hệ thống.\\ \midrule
% \textbf{Hậu điều kiện}     & Tạo phiên đăng nhập hợp lệ. \\ \midrule
% \textbf{Luồng sự kiện chính} & 
% \begin{tabular}[t]{@{}p{10.5cm}@{}}
% 1. Chọn phương thức đăng nhập bằng Google.\\
% 2. Hệ thống điều hướng sang trang đăng nhập với Google.\\
% 3. Người dùng chọn account muốn đăng nhập.\\
% 4. Hệ thống xác thực.\\
% 5. Hệ thống tạo session/JWT và chuyển vào ứng dụng. 
% \end{tabular} \\ \midrule
% \textbf{Luồng ngoại lệ} &
% \begin{tabular}[t]{@{}p{10.5cm}@{}}
% N/A
% \end{tabular} \\
% \bottomrule
% \end{tabular}
% \caption{Mô tả Use Case: Đăng nhập tài khoản Google}
% \label{usecase-8-02}
% \end{table}

% \begin{table}[H]
% \centering
% \renewcommand{\arraystretch}{1.2}
% \begin{tabular}{@{}p{4cm} p{11cm}@{}}
% \toprule
% \textbf{Use case ID}       & UC-M8-05 \\ \midrule
% \textbf{Tên Use case}      & Đăng nhập bằng email/sđt \\ \midrule
% \textbf{Các tác nhân}      & Người dùng chung \\ \midrule
% \textbf{Mô tả}             & Cho phép người dùng đăng nhập vào hệ thống bằng tài khoản email/sđt đã đăng ký\\ \midrule
% \textbf{Trigger}    & Người dùng bấm “Đăng nhập".\\ \midrule
% \textbf{Tiền điều kiện}    &  Người dùng đã có tài khoản trong hệ thống.\\ \midrule
% \textbf{Hậu điều kiện}     & Tạo phiên đăng nhập hợp lệ. \\ \midrule
% \textbf{Luồng sự kiện chính} & 
% \begin{tabular}[t]{@{}p{10.5cm}@{}}
% 1. Chọn phương thức đăng nhập bằng Email/SDT.\\
% 2. Hệ thống hiển thị form đăng nhập.\\
% 3. Người dùng điền email/sdt và mật khẩu.\\
% 4. Hệ thống xác thực tài khoản.\\
% 5. Hệ thống tạo session/JWT và chuyển vào ứng dụng.
% \end{tabular} \\ \midrule
% \textbf{Luồng ngoại lệ} &
% \begin{tabular}[t]{@{}p{10.5cm}@{}}
% N/A
% \end{tabular} \\
% \bottomrule
% \end{tabular}
% \caption{Mô tả Use Case: Đăng nhập tài khoản Email/SĐT}
% \label{usecase-8-05}
% \end{table}

% \begin{table}[H]
% \centering
% \renewcommand{\arraystretch}{1.2}
% \begin{tabular}{@{}p{4cm} p{11cm}@{}}
% \toprule
% \textbf{Use case ID}       & UC-M8-07 \\ \midrule
% \textbf{Tên Use case}      & Xem danh sách tài khoản vi phạm \\ \midrule
% \textbf{Các tác nhân}      & Quản trị viên hệ thống \\ \midrule
% \textbf{Mô tả}             & Xem danh sách các tài khoản bị báo cáo vi phạm quy định..\\ \midrule
% \textbf{Trigger}    & Quản trị viên mở trang “Tài khoản vi phạm”.\\ \midrule
% \textbf{Tiền điều kiện}    &  Người dùng đăng nhập với vai trò quản trị viên hệ thống\\ \midrule
% \textbf{Hậu điều kiện}     & Hiển thị danh sách tài khoản vi phạm. \\ \midrule
% \textbf{Luồng sự kiện chính} & 
% \begin{tabular}[t]{@{}p{10.5cm}@{}}
% 1. Hệ thống hiển thị danh sách (phân trang).\\
% 2. Lọc theo mức độ, thời gian, loại vi phạm.\\
% \end{tabular} \\ \midrule
% \textbf{Luồng ngoại lệ} &
% \begin{tabular}[t]{@{}p{10.5cm}@{}}
% N/A
% \end{tabular} \\
% \bottomrule
% \end{tabular}
% \caption{Mô tả Use Case: Xem danh sách tài khoản vi phạm}
% \label{usecase-8-07}
% \end{table}

% \begin{table}[H]
% \centering
% \renewcommand{\arraystretch}{1.2}
% \begin{tabular}{@{}p{4cm} p{11cm}@{}}
% \toprule
% \textbf{Use case ID}       & UC-M8-09 \\ \midrule
% \textbf{Tên Use case}      & Khóa / Mở tài khoản \\ \midrule
% \textbf{Các tác nhân}      & Quản trị viên hệ thống \\ \midrule
% \textbf{Mô tả}             & Xử lý tài khoản vi phạm bằng cách thay đổi trạng thái tài khoản.\\ \midrule
% \textbf{Trigger}    & Quản trị viên mở trang “Tài khoản vi phạm”.\\ \midrule
% \textbf{Tiền điều kiện}    &  Người dùng đăng nhập với vai trò quản trị viên hệ thống \& trạng thái hiện tại cho phép thao tác (đang mở → khóa; đang khóa → mở). \\ \midrule
% \textbf{Hậu điều kiện}     &  Trạng thái tài khoản thay đổi tương ứng; thông báo được gửi. \\ \midrule
% \textbf{Luồng sự kiện chính} & 
% \begin{tabular}[t]{@{}p{10.5cm}@{}}
% 1. QTV chọn tài khoản và hành động Khóa hoặc Mở.
% 2. QTV nhập/chọn lý do và xác nhận thao tác.
% 3. Hệ thống cập nhật trạng thái.
% \end{tabular} \\ \midrule
% \textbf{Luồng ngoại lệ} &
% \begin{tabular}[t]{@{}p{10.5cm}@{}}
% Lỗi hệ thống → hủy thao tác, giữ nguyên trạng thái.
% \end{tabular} \\
% \bottomrule
% \end{tabular}
% \caption{Mô tả Use Case: Khóa \/ Mở tài khoản}
% \label{usecase-8-09}
% \end{table}
